% Lösungen Einheit 3 von 3 – Maximum bestimmen und deuten (zusammenbau v0.1)
\einheitenkopf{Einheit 3 von 3 · Maximum bestimmen und deuten}

\erg{13}{a) der Flächeninhalt \quad b) $A'(a) = 30 - 4a = 0$ liefert $a = 7{,}5$ (im Definitionsbereich) \quad c) $A'(a) = 24 - 4a = 0$ liefert $a = 6$; $A''(6) = -4 < 0$, also Maximum; $b = 24 - 12 = 12$ m, $a = 6$ m, $A(6) = 72$ m² \quad d) $A'(x) = -3x^4 + 6x^2 + 24 = 0$; mit $z = x^2$: $z^2 - 2z - 8 = 0$, $z_1 = 4$, $z_2 = -2$; $z_2$ liefert keine Lösung, aus $z_1$ folgt $x = 2$ ($x = -2$ liegt nicht im Intervall); $A''(2) = -72 < 0$ – Maximum \quad e) $A(u) = u \cdot f(u) = (u^2 + 3u) \cdot e^{-0{,}5u}$; $A'(u) = 0$ liefert $u^2 - u - 6 = 0$, also $u = 3$ ($u = -2$ entfällt); $P(3 | 1{,}34)$ \quad f) $d(x) = f(x) - g(x) = -0{,}01x^2 + 0{,}5x$; $d'(x) = -0{,}02x + 0{,}5 = 0$ liefert $x = 25$; $d(25) = -6{,}25 + 12{,}5 = 6{,}25$ m \quad g) $T(u) = (10 - u) \cdot (10u - u^2) = u \cdot (10 - u)^2$; $T'(u) = (10 - u) \cdot (10 - 3u) = 0$ liefert $u = 10$ (nicht im Bereich) oder $u = \tfrac{10}{3} \approx 3{,}33$; Vorzeichenwechsel von $+$ nach $-$ – Maximum \quad h) $A(x) = x \cdot f(x) = 4x - x^3$, $A'(1) = 4 - 3 = 1 \ne 0$ – die notwendige Bedingung ist verletzt, also ist $x = 1$ keine Maximalstelle. \quad i) Aufgabe: Die Punkte $(-3 | 0)$, $(u | 0)$ und $(u | f(u))$ mit $-3 < u < 3$ bilden ein rechtwinkliges Dreieck; bestimme $u$ so, dass sein Flächeninhalt maximal wird, und gib diesen an. (1) Zielfunktion aus Grundseite $u + 3$ und Höhe $f(u)$; (2) notwendige Bedingung; (3) hinreichende Bedingung – Maximum; (4) der größte Flächeninhalt ist $A(1) = 16$.}
\erg{14}{a) $d(x)^2 = (x - 2)^2 + (\sqrt{x})^2 = x^2 - 3x + 4$; Ableitung $2x - 3 = 0$ liefert $x = 1{,}5$; kleinster Abstand $\sqrt{1{,}75} \approx 1{,}32$}
\erg{15}{a) $x = -5$ liegt nicht im Definitionsbereich und gehört zu keinem Rechteck; das einzige größte Rechteck gehört zu $x = 5$ mit $A(5) = 250$. \quad b) Für $x$ gegen 0 schrumpft die Breite, für $x$ gegen 2 die Höhe gegen null; beide Randfälle sind ausgeschlossen, die Fläche wird beliebig klein, erreicht aber nie null – ein kleinstes gibt es nicht; dazwischen nimmt die Flächenfunktion einen größten Wert an. \quad c) $V'(x) = 12x^2 - 240x + 900 = 0$ liefert $x = 5$ oder $x = 15$ (Rand, keine Schachtel); $V''(5) = -120 < 0$; Quadrate mit 5 cm Seitenlänge, Volumen $V(5) = 2\,000$ cm³}
